% ============================================================
\section{实验验证}
% ============================================================
\sectioncover{03}{实验验证}{}

\begin{frame}{实验设置}
  \slidekey{在此填写本页核心结论。}

  \begin{block}{模型与训练配置}
    \begin{compactitem}
      \item 主实验模型：模型名称
      \item 跨模型验证：模型名称
      \item 关键超参数一
      \item 关键超参数二
      \item 关键超参数三
    \end{compactitem}
  \end{block}

  \vspace{0.10cm}
  \begin{exampleblock}{方法配置}
    \begin{compactitem}
      \item 配置项一
      \item 配置项二
      \item 配置项三
    \end{compactitem}
  \end{exampleblock}
\end{frame}

% ====== 图片占位 ======

\begin{frame}{数据分布}
  \slidekey{在此填写本页核心结论。}

  \centering
  % \includegraphics[width=0.80\paperwidth,height=0.52\paperheight,keepaspectratio]{distribution.pdf}
  \vspace{3cm}
  \textcolor{ZJUGray}{\large [ 数据分布图占位 ]}

  \vspace{0.12cm}
  \begin{block}{}
    \small
    \begin{compactitem}
      \item 数据分布说明一。
      \item 数据分布说明二。
    \end{compactitem}
  \end{block}
\end{frame}

\begin{frame}{主实验结果图}
  \slidekey{在此填写本页核心结论。}

  \centering
  % \includegraphics[width=0.88\paperwidth,height=0.65\paperheight,keepaspectratio]{results_chart.pdf}
  \vspace{4cm}
  \textcolor{ZJUGray}{\large [ 主实验结果图占位 ]}
\end{frame}

% ====== 保留一个空表格结构 ======

\begin{frame}{实验结果数据}
  \slidekey{在此填写本页核心结论。}

  \centering
  \resizebox{\textwidth}{!}{%
    \smalltab
    \begin{tabular}{lccccccc}
      \toprule
      \textbf{模型} & \textbf{指标一} & \textbf{指标二} & \textbf{指标三} & \textbf{指标四} & \textbf{指标五} & \textbf{指标六} & \textbf{指标七} \\
      \midrule
      基线模型 & & & & & & & \\
      \rowcolor{SOTAPink}
      \textbf{本文方法} & & & & & & & \\
      $\\Delta$ & & & & & & & \\
      \bottomrule
    \end{tabular}
  }

  \vspace{0.28cm}
  \tagbox[ZJUGreen]{在此填写实验结果的总结性评语}
\end{frame}

\begin{frame}{人工评测}
  \slidekey{在此填写本页核心结论。}

  \centering
  \small 评测说明文字。

  \vspace{0.20cm}
  \begin{tabular}{lccc}
    \toprule
    \textbf{对比} & \textbf{本文方法} & \textbf{基线} & \textbf{平局} \\
    \midrule
    设置一 & & & \\
    设置二 & & & \\
    \bottomrule
  \end{tabular}
\end{frame}

\begin{frame}{消融实验}
  \slidekey{在此填写本页核心结论。}

  \centering
  \begin{tabular}{cccc}
    \toprule
    \textbf{配置} & \textbf{指标一} & \textbf{指标二} & \textbf{指标三} \\
    \midrule
    基线 & & & \\
    变体一 & & & \\
    \rowcolor{SOTAPink}
    \textbf{最优} & & & \\
    变体二 & & & \\
    \bottomrule
  \end{tabular}

  \vspace{0.20cm}
  \begin{tikzpicture}
    \node[
      fill=TitleBlue,
      draw=ZJUBlue!35,
      rounded corners=3pt,
      text width=0.90\textwidth,
      align=center,
      inner sep=0.14cm,
      font=\small
    ] {消融实验的结论总结。};
  \end{tikzpicture}
\end{frame}

% ====== 图片占位 ======

\begin{frame}{定性分析}
  \slidekey{在此填写本页核心结论。}

  \centering
  % \includegraphics[width=0.84\paperwidth,height=0.56\paperheight,keepaspectratio]{qualitative.pdf}
  \vspace{4cm}
  \textcolor{ZJUGray}{\large [ 定性分析图占位 ]}

  \vspace{0.14cm}
  \begin{tikzpicture}
    \node[
      fill=TitleBlue,
      draw=ZJUBlue!35,
      rounded corners=3pt,
      text width=0.90\textwidth,
      align=center,
      inner sep=0.14cm,
      font=\small
    ] {\textbf{基线}：基线结果描述。\\[2pt]\textbf{本文}：本文方法结果描述。};
  \end{tikzpicture}
\end{frame}

\begin{frame}{训练过程分析}
  \slidekey{在此填写本页核心结论。}

  \centering
  % \includegraphics[width=0.80\paperwidth,height=0.55\paperheight,keepaspectratio]{training_curve.pdf}
  \vspace{4cm}
  \textcolor{ZJUGray}{\large [ 训练曲线图占位 ]}

  \vspace{0.12cm}
  \small 训练过程的分析说明。
\end{frame}
